High-resolution cascaded MIMO radar produces more raw data than a vehicle's
network and central processor can comfortably absorb. Compressing that data at
the sensor buys back bandwidth, but only if the compressor is lossless, fast and
small, because anything it discards is lost before target detection ever sees
it.

This thesis compares data compression methods for automotive radar on real data
and implements the most suitable of them in hardware. Eight methods were
evaluated in MATLAB on 50 frames of the public ColoRadar dataset, recorded with
a twelve-transmitter Texas Instruments cascade radar. The lossy quantisers
reached compression ratios of up to \num{3.76} but missed up to
\SI{28.7}{\percent} of targets, and floating-point storage barely compressed at
all. Two lossless predictive coders compressed by more than three:
Doppler Redundancy Hybrid Encoding (DRHE), which predicts magnitude and phase
with fixed filters, and a second-order linear predictive coder (LPC) whose
coefficients are fitted to each frame. On a time-division MIMO radar both must
predict each ramp from the previous ramp of the \emph{same transmitter}, twelve
ramps back; predicting from the adjacent ramp does worse than not predicting
at all.

The compression ratio was broken down into its sources. Huffman coding of the
16-bit range-FFT data alone achieves \num{3.48}; prediction adds a further
\num{0.22}--\num{0.27}; and a Huffman dictionary re-derived for the prediction
errors adds \num{0.47} for DRHE and \num{0.71} for LPC, giving \num{4.22} and
\num{4.40} overall, measured on the hardware, with every frame reconstructed
exactly.

Both compressors and their decompressors were written in C++, synthesised
with Vitis HLS and Vivado, and run on an AMD KCU116 (Kintex UltraScale+) board
at \SI{100}{\mega\hertz}. On the board all 50 frames were compressed
byte-identically to the simulation, and the board's own decoders reconstructed
every frame exactly. DRHE
compresses a frame in \SI{0.61}{\milli\second} (\SI{5.1}{Gbit/s}) and streams
its output with \SI{13}{\micro\second} of latency. LPC takes
\SI{3.25}{\milli\second} and, because its coefficients depend on the whole
frame, cannot output anything until the whole frame has arrived: its last
output word follows its last input word by \SI{3.0}{\milli\second}. The
decoders each reconstruct one value per clock, about \SI{2}{\milli\second} per
frame. Encoder and decoder together use about \SI{23}{\percent} of the chip's
logic for either algorithm. LPC
compresses about \SI{4}{\percent} better; DRHE is recommended for a streaming
sensor front end, because it needs no frame buffer and has a small fraction of
LPC's latency.
